\subsection*{小结}

\begin{table}[htbp]\centering\small
\begin{tabularx}{\linewidth}{@{}lXX@{}}
\toprule
\textbf{概念} & \textbf{原语} & \textbf{何时使用} \\
\midrule
谓词执行 & \co@if@ & 数据相关的判断（边界、条件） \\
结构化并发 & \co@inthreads.async@ & 编译时划分线程（生产者／消费者、异构角色） \\
区域间通信 & \co@shared event@ / \co@wait@ / \co@trigger@ & 并发区域之间的协调 \\
主机并发 & \co@parallel.async@ / \co@stream s@ & 多核函数重叠、非阻塞启动 \\
持久调度 & \co@if@ + \co@foreach@ + \co@#@ & 固定网格规模、带填充守卫的分块条带化 \\
\bottomrule
\end{tabularx}
\end{table}

完整流水线为前面的骨架补上两项依赖：计算前输入必须就绪，复用槽位前操作数必须用完。修改流水线时，可以沿一个槽位追踪完整周期，再追踪下一次复用。与只看某一侧循环相比，这样更容易发现缺失信号和迭代次数不匹配。
